\section{Introduction}
% Present the problem, motivation, research gap, proposed solution, research questions, and main contributions

Weather radar systems play an important role in monitoring precipitation, characterizing severe meteorological phenomena, and supporting short-term weather forecasting and early-warning applications. Radar reflectivity fields provide spatial information about the distribution and intensity of hydrometeors, allowing the structure and evolution of precipitation systems to be analyzed. However, radar observations are affected by physical and acquisition-related limitations, including clutter, attenuation, beam blockage, and range-dependent reductions in spatial resolution, which can impair the representation of fine-scale meteorological structures \cite{green2024assessing}. Such limitations are particularly relevant for localized convective cells, intense echoes, and sharp precipitation boundaries, whose spatial characteristics may be difficult to preserve under degraded or low-resolution observations \cite{li2026fine}. Consequently, radar-image restoration and super-resolution have gained attention as computational approaches for improving the spatial fidelity of reflectivity fields and reconstructing finer meteorological structures from degraded observations \cite{shi2023radar}.


Deep learning has significantly advanced weather-radar image restoration and super resolution. CNN-based architectures learn nonlinear mappings between low- and high-resolution radar fields, while generative adversarial approaches have been investigated to recover finer details and stronger intensity gradients \cite{xu2025generalized}. More recently, diffusion-based generative approaches have been explored for radar-data enhancement and reconstruction \cite{peng2026multi}. Residual attention mechanisms and multi-scale feature extraction have also been incorporated to better capture spatial dependencies and reconstruct detailed radar structures \cite{yu2023weather}. Despite these advances, restoration methods are generally developed and assessed as individual models, with their effectiveness commonly characterized through aggregate performance metrics over entire test sets \cite{liu2025msrgan}. Such aggregate evaluation
 provides limited insight into whether different architectures exhibit complementary strengths under distinct radar conditions, motivating the investigation of restoration models according to their performance across heterogeneous radar-image regimes rather than solely according to overall reconstruction quality.

 Weather-radar imagery exhibits substantial heterogeneity in the spatial distribution, morphology, and intensity of precipitation echoes, creating restoration requirements that may vary considerably across image conditions. Precipitation fields can contain markedly different structures, ranging from widespread stratiform regions to compact convective cells with stronger reflectivity and pronounced spatial variability \cite{mikidadi2025application}. Beyond these meteorological differences, restoration performance may also vary with factors such as echo intensity, spatial structure, location within the radar domain, and degradation severity. Consequently, a model that achieves strong aggregate reconstruction performance is not necessarily expected to be equally effective under every radar-image condition. This variability motivates the use of heterogeneous ensembles, in which models with different architectural characteristics and reconstruction behaviours can provide complementary capabilities across distinct radar-image regimes \cite{sun2024ensir}. Rather than relying solely on a single globally best model, such complementarity provides a basis for selecting models according to the regimes in which they demonstrate sufficient restoration capability.


 Although a heterogeneous ensemble can exploit complementary model capabilities, retaining the full set of models may incur considerable computational overhead in terms of inference latency, memory footprint, and processing cost. Ensemble pruning provides a means to select a compact subset of models while preserving the benefits of the original ensemble and reducing its computational burden. Existing studies have approached ensemble pruning through performance-based ranking, diversity driven selection, multi-objective optimisation, and mathematical-programming formulations \cite{rezk2024metaheuristic}. In particular, quadratic and integer programming approaches have been employed to balance predictive performance, ensemble diversity, and subset size \cite{emine2025free}, while Pareto-based methods explicitly explore trade-offs among competing objectives \cite{he2024margin}, \cite{qiao2024ensemble}. These studies establish that optimisation-based ensemble pruning is not itself a new concept; rather, a central question is which properties of the candidate models should determine their inclusion in the selected subset. Current criteria remain predominantly global relying on aggregate reconstruction metrics or generic diversity measures and therefore do not explicitly guarantee coverage of the heterogeneous meteorological and observational regimes present in weather-radar imagery.

 In this paper, we propose a cost-aware Weighted Set Multicover formulation for pruning a heterogeneous ensemble of deep-learning models for weather-radar image restoration. Rather than selecting models solely according to aggregate reconstruction performance or generic diversity, each candidate model is evaluated according to its ability to satisfy predefined performance requirements across distinct radar-image regimes. Each regime is represented as an element requiring coverage, and a model covers a regime when it meets specified criteria related to reconstruction quality and the preservation of meteorologically relevant structures. The multicover formulation further allows important or rare precipitation regimes to require coverage by multiple competent models. The optimization objective accounts for both the number and computational cost of the selected models, seeking a compact subset while maintaining sufficient coverage of the radar-image regimes. The proposed approach therefore integrates regime-oriented coverage, preservation of critical meteorological structures, and computational efficiency within a unified ensemble-pruning formulation. Its effectiveness is evaluated against the best individual model, the complete ensemble, random and Top-k selection, diversity-based pruning, Pareto-based pruning, and integer-programming-based ensemble selection.